In this work, we explore generating Wikipedia-like articles from scratch as a way to push the frontier of automatic expository writing and long-form article generation. While our approach significantly outperforms baseline methods in both automatic and human evaluations, the quality of machine-written articles still lags behind well-revised human-authored articles, specifically in aspects of neutrality and verifiability. Although \system discovers different perspectives in researching the given topic, the collected information may still be biased towards dominant sources on the Internet and may contain promotional content. Moreover, the verifiability issues identified in this work go beyond factual hallucination, which highlights new challenges to grounded writing systems.

Another limitation of this work is that although we focus on the task of generating Wikipedia-like articles from scratch, our task setup is still simplified to only consider the generation of free-form text. Human-authored high-quality Wikipedia articles usually contain structured data and multi-modal information. We leave the exploration of generating multi-modal grounded articles for future work.
